\documentclass[tikz,border=6pt]{standalone}
\usepackage{peabra}
\begin{document}
\begin{tikzpicture}[x=1cm,y=1cm,
  ism/.style={doboz=14mm, font=\footnotesize, minimum height=9mm},
  ker/.style={kiemeltdoboz=14mm, font=\footnotesize\bfseries, minimum height=9mm},
  rsz/.style={sotetdoboz=17mm, font=\footnotesize\bfseries, minimum height=11mm},
  rszk/.style={kiemeltdoboz=17mm, font=\footnotesize\bfseries, minimum height=11mm}]
  \foreach \x/\cim/\al/\b/\r/\k in {
      0/{Rendszerelemzés}/{direkt feladat}/ism/rsz/ker,
      6.2/{Rendszertervezés}/{1. inverz feladat}/ism/rszk/ism,
      12.4/{Rendszerirányítás}/{2. inverz feladat}/ker/rsz/ism} {
    \begin{scope}[xshift=\x cm]
      \node[\b] (b) at (-2.15,0) {bemenet};
      \node[\r] (r) at (0,0) {rendszer};
      \node[\k] (k) at (2.15,0) {kimenet};
      \draw[nyil] (b) -- (r); \draw[nyil] (r) -- (k);
      \node[font=\cimbetu\small, text=peSotet] at (0,1.25) {\cim};
      \node[font=\scriptsize, text=peSzurke] at (0,0.8) {\al};
    \end{scope}}
  \node[font=\scriptsize, text=peSotet, align=center, text width=52mm] at (0,-1.15) {adott a bemenet és a rendszer;\\keressük a kimenetet\\{\color{peSzurke}(pl. mennyit termel a cella?)}};
  \node[font=\scriptsize, text=peSotet, align=center, text width=52mm] at (6.2,-1.15) {adott a bemenet és a kívánt kimenet;\\keressük a rendszert\\{\color{peSzurke}(milyen cella termel 140 ezer darabot?)}};
  \node[font=\scriptsize, text=peSotet, align=center, text width=52mm] at (12.4,-1.15) {adott a rendszer és a kívánt kimenet;\\keressük a bemenetet\\{\color{peSzurke}(hogyan ütemezzük a rendeléseket?)}};
  \node[font=\scriptsize, text=peKiemel] at (6.2,-2.25) {narancs: a keresett elem};
\end{tikzpicture}
\end{document}
